\section{Write a canonical container around the selected PCM}
The output writer uses a twelve-byte RIFF prefix, a twenty-four-byte format chunk
with sixteen payload bytes, and an eight-byte data header followed by PCM. Its PCM
offset is always 44. It derives the sizes from the validated rate, channel count and
frame count in the description prefix; it never copies the parsed input offset 56
into an output pointer calculation.

For the worked four-frame stereo result, the PCM size is sixteen and the complete
canonical file size is sixty. The RIFF size field is therefore 52. Keep rate 16000,
channel count two, byte rate 64000, frame size four and width sixteen. Only the PCM
codes change during filtering. The input's JUNK and note chunks are not carried into
this canonical output. Copying PCM without filtering would preserve its audio bytes, but would not
preserve the original container or its ancillary metadata.

Validate the description and require the entire $44+2CN$ output capacity before
reading PCM or writing a header. On description or capacity failure, leave output
unchanged. On success, copy exactly the selected PCM and return the complete file
length, with no trailing terminator. The zero-frame case writes a valid forty-four-byte
file and reads no PCM bytes. The writer has no reason to emit an odd data pad because
every selected frame contains an even number of bytes.
